\chapter{Primes}

\begin{center}\itshape
There is always a prime between $n$ and $2n$. Chebyshev proved it in 1852;\\
Erd\H{o}s, at nineteen, found the proof given here, which uses nothing but the binomial coefficient $\binom{2n}n$.
\end{center}

\section{Infinitely many primes}

\begin{theorem}[Euclid]
There are infinitely many primes.
\end{theorem}
\begin{proof}
Given primes $p_1,\dots,p_k$, any prime factor of $p_1\cdots p_k+1$ is none of them.
\end{proof}

Euler went much further: the primes are numerous enough for their reciprocals to diverge, whereas $\sum\frac1{n^2}$ converges. So primes are, in a precise sense, denser than squares.

\begin{theorem}[Euler]\label{thm:euler}
For $x\ge2$, $\displaystyle\sum_{p\le x}\frac1p\ge\ln\ln x-1$. In particular $\sum_p\frac1p=+\infty$.
\end{theorem}
\begin{proof}
Every $n\le x$ factors into primes $\le x$, so expanding the finite product of geometric series,
\[
\prod_{p\le x}\left(1-\frac1p\right)^{-1}=\prod_{p\le x}\sum_{k\ge0}\frac1{p^k}\ge\sum_{n\le x}\frac1n\ge\int_1^{\lfloor x\rfloor+1}\frac{\mathrm dt}t\ge\ln x.
\]
For $0\le t<1$, $-\ln(1-t)=\sum_{k\ge1}\frac{t^k}k\le t+\sum_{k\ge2}t^k=t+\frac{t^2}{1-t}$. With $t=\frac1p$ the last term is $\frac1{p(p-1)}$, and $\sum_p\frac1{p(p-1)}\le\sum_{n\ge2}\frac1{n(n-1)}=1$. Taking logarithms,
$\ln\ln x\le\sum_{p\le x}-\ln\left(1-\frac1p\right)\le\sum_{p\le x}\frac1p+1$.
\end{proof}

\section{Chebyshev's bounds}

Let $\pi(x)$ be the number of primes $\le x$. The prime number theorem, proved in Tome~5, says $\pi(x)\sim\frac x{\ln x}$. Chebyshev's elementary bounds already give the right order of magnitude. Everything rests on the central binomial coefficient.

\begin{lemma}\label{lem:binom}
For $n\ge1$:
\begin{enumerate}[label=(\roman*)]
\item $\dfrac{4^n}{2n}\le\dbinom{2n}n\le4^n$, and $\dbinom{2n+1}n\le4^n$;
\item every prime power $p^a$ dividing $\binom{2n}n$ exactly satisfies $p^a\le2n$;
\item primes $p$ with $\sqrt{2n}<p$ divide $\binom{2n}n$ at most once, and primes with $\frac{2n}3<p\le n$ do not divide it at all, if $n\ge5$.
\end{enumerate}
\end{lemma}
\begin{proof}
(i) $\binom{2n}n$ is the largest of the $2n+1$ terms of $\sum_k\binom{2n}k=4^n$, and is at least the sum $2$ of the two end terms; so $4^n\le(2n-1)\binom{2n}n+2\le2n\binom{2n}n$. Also $\binom{2n+1}n=\binom{2n+1}{n+1}$ and both appear in $\sum_k\binom{2n+1}k=2\cdot4^n$.

(ii) By Legendre (Chapter~1), $v_p\binom{2n}n=\sum_{k\ge1}\left(\left\lfloor\frac{2n}{p^k}\right\rfloor-2\left\lfloor\frac n{p^k}\right\rfloor\right)$. Each term is $0$ or $1$ (as $\lfloor2y\rfloor-2\lfloor y\rfloor\in\{0,1\}$), and is $0$ when $p^k>2n$. So $v_p\le\max\{k:p^k\le2n\}$.

(iii) If $p>\sqrt{2n}$ only the term $k=1$ can be non-zero. If $\frac{2n}3<p\le n$ then $2\le\frac{2n}p<3$ and $1\le\frac np<\frac32$, so the $k=1$ term is $2-2=0$, and $p^2>2n$ when $n\ge5$ (as then $\frac{2n}3>\sqrt{2n}$).
\end{proof}

\begin{theorem}[Chebyshev-type bounds]\label{thm:cheb}
\begin{enumerate}[label=(\roman*)]
\item For every real $x\ge1$, $\displaystyle\prod_{p\le x}p<4^x$.
\item For every $n\ge1$, $\pi(2n)\ge\dfrac{n\ln4-\ln(2n)}{\ln(2n)}$. In particular $\pi(x)\ge\frac{\ln2}2\cdot\frac x{\ln x}$ for large $x$.
\end{enumerate}
\end{theorem}
\begin{proof}
(i) It suffices to treat integers $x=N$, by induction. $N=1,2$ are clear, and for even $N>2$ the product equals the one for $N-1$. For $N=2m+1$, every prime in $]m+1,2m+1]$ divides $\binom{2m+1}m=\frac{(2m+1)!}{m!\,(m+1)!}$ (it divides the numerator, not the denominator), so their product is at most $\binom{2m+1}m\le4^m$, and
$\prod_{p\le2m+1}p\le4^m\prod_{p\le m+1}p<4^m\cdot4^{m+1}=4^{2m+1}$.

(ii) Only primes $p\le2n$ divide $\binom{2n}n$, each to a power at most $2n$ by the lemma, so $\frac{4^n}{2n}\le\binom{2n}n\le(2n)^{\pi(2n)}$. Take logarithms.
\end{proof}

\section{Bertrand's postulate}

\begin{theorem}[Bertrand's postulate; Chebyshev 1852, proof by Erd\H{o}s 1932]\label{thm:bertrand}
For every $n\ge1$ there is a prime $p$ with $n<p\le2n$.
\end{theorem}
\begin{proof}
\emph{Small $n$.} The primes $2,3,5,7,13,23,43,83,163,317,631$ each are less than twice the previous one. For $1\le n<631$, let $q$ be the largest of them with $q\le n$ (or $q=1$ if $n=1$) and $q'$ the next one: then $n<q'<2q\le2n$. So the theorem holds for $n\le630$.

\emph{Large $n$.} Let $n\ge468$ and suppose there is no prime in $]n,2n]$. By Lemma~\ref{lem:binom}, the primes dividing $\binom{2n}n$ are all $\le\frac{2n}3$. Those $\le\sqrt{2n}$ (fewer than $\sqrt{2n}$ of them) contribute at most $2n$ each; the others contribute at most $p$, and their product is at most $\prod_{p\le2n/3}p<4^{2n/3}$ by Theorem~\ref{thm:cheb}. Hence
\[
\frac{4^n}{2n}\le\binom{2n}n<(2n)^{\sqrt{2n}}\,4^{2n/3},\qquad\text{so}\qquad4^{n/3}<(2n)^{1+\sqrt{2n}}.
\]
Taking logarithms, $\frac n3\ln4<\left(1+\sqrt{2n}\right)\ln(2n)$, which is false for every $n\ge468$ (the left side grows like $n$, the right like $\sqrt n\ln n$; the companion notebook checks the crossover). This contradiction proves the theorem.
\end{proof}

\section{Exercises}

\begin{exercise}[$\star$]\label{ex:mod4}
Show that there are infinitely many primes $p\equiv3\pmod4$. \emph{Hint: consider $4p_1\cdots p_k-1$.}
\end{exercise}

\begin{exercise}[$\star$]\label{ex:mersenne}
Show that if $2^n-1$ is prime then $n$ is prime, and that if $2^n+1$ is prime ($n\ge1$) then $n$ is a power of $2$.
\end{exercise}

\begin{exercise}[$\star$]\label{ex:germain}
Show that $n^4+4$ is composite for every integer $n>1$.
\end{exercise}

\begin{exercise}[$\star\star$]\label{ex:fermatnum}
Let $\mathcal F_n=2^{2^n}+1$ be the Fermat numbers. Show that $\mathcal F_0\mathcal F_1\cdots\mathcal F_{n-1}=\mathcal F_n-2$, deduce that the Fermat numbers are pairwise coprime, and give a new proof that there are infinitely many primes.
\end{exercise}

\begin{exercise}[$\star\star$]\label{ex:gaps}
Show that there are arbitrarily long runs of consecutive composite integers. Why does this not contradict Bertrand's postulate?
\end{exercise}

\begin{exercise}[$\star\star$]\label{ex:factorial}
Show that $n!$ is not a perfect power ($k$-th power with $k\ge2$) for $n\ge2$.
\end{exercise}

\begin{exercise}[$\star\star\star$]\label{ex:erdosdiv}
\emph{Erd\H{o}s' proof that $\sum\frac1p$ diverges.} Let $p_1<p_2<\cdots$ be the primes, and suppose $\sum_{j>k}\frac1{p_j}<\frac12$ for some $k$. Let $N\ge1$.
\begin{enumerate}[label=(\alph*)]
\item Show that at most $\sum_{j>k}\frac N{p_j}<\frac N2$ integers $n\le N$ have a prime factor $p_j$ with $j>k$.
\item Writing $n=ab^2$ with $a$ squarefree, show that at most $2^k\sqrt N$ integers $n\le N$ have all their prime factors among $p_1,\dots,p_k$.
\item Conclude by taking $N=2^{2k+2}$.
\end{enumerate}
\end{exercise}

\begin{exercise}[$\star\star\star$]\label{ex:furstenberg}
\emph{Furstenberg's topological proof (1955).} Call $U\subset\Z$ open if it is a union of arithmetic progressions $a+b\Z$ with $b\ge1$.
\begin{enumerate}[label=(\alph*)]
\item Show that this defines a topology on $\Z$ in which every non-empty open set is infinite.
\item Show that each $a+b\Z$ is also closed.
\item If there were finitely many primes, show that $\{-1,1\}$ would be open, and conclude.
\end{enumerate}
\end{exercise}

\begin{beauty}
Euclid gives infinitely many primes, Euler shows they are plentiful ($\sum\frac1p$ diverges like $\ln\ln x$), and Chebyshev pins their count to the order $\frac x{\ln x}$. The exact asymptotic, the prime number theorem, needs complex analysis and is the summit of Tome~5. Primes in arithmetic progressions (Exercise~\ref{ex:mod4} is the easy case) need Dirichlet's characters, built in Tome~2. And Exercise~\ref{ex:furstenberg} shows that even topology has something to say about primes.
\end{beauty}

\section{Solutions}

\begin{solution}{ex:mod4}
Suppose $p_1,\dots,p_k$ are all the primes $\equiv3\pmod4$, and let $N=4p_1\cdots p_k-1\equiv3\pmod4$. $N$ is odd, and a product of primes all $\equiv1\pmod4$ is $\equiv1\pmod4$, so $N$ has a prime factor $q\equiv3\pmod4$. But no $p_i$ divides $N$ (it would divide $1$): contradiction.
\end{solution}

\begin{solution}{ex:mersenne}
If $n=ab$ with $1<a<n$, then $2^a-1$ divides $2^{ab}-1=(2^a)^b-1$ (as $x-1\mid x^b-1$) and $1<2^a-1<2^n-1$. If $n$ has an odd factor $m>1$, $n=mk$, then $2^k+1$ divides $2^n+1=(2^k)^m+1$ (as $x+1\mid x^m+1$ for $m$ odd), and $1<2^k+1<2^n+1$.
\end{solution}

\begin{solution}{ex:germain}
$n^4+4=(n^2+2)^2-(2n)^2=(n^2+2n+2)(n^2-2n+2)$, and $n^2-2n+2=(n-1)^2+1\ge2$ for $n>1$, while the other factor is larger.
\end{solution}

\begin{solution}{ex:fermatnum}
By induction: true for $n=1$ ($\mathcal F_0=3=5-2$), and $\left(\mathcal F_n-2\right)\mathcal F_n=\left(2^{2^n}-1\right)\left(2^{2^n}+1\right)=2^{2^{n+1}}-1=\mathcal F_{n+1}-2$. A common divisor $d$ of $\mathcal F_m$ and $\mathcal F_n$ with $m<n$ divides $\mathcal F_n-\mathcal F_0\cdots\mathcal F_{n-1}=2$, and Fermat numbers are odd, so $d=1$. Choosing a prime factor of each $\mathcal F_n$ gives infinitely many distinct primes.
\end{solution}

\begin{solution}{ex:gaps}
For $2\le j\le n+1$, $j$ divides $(n+1)!+j$, so the $n$ consecutive integers $(n+1)!+2,\dots,(n+1)!+(n+1)$ are composite. No contradiction: these gaps start around $(n+1)!$, far beyond $n$, and Bertrand only forbids a gap from $m$ to $2m$; a gap of length $n$ starting at $m\approx(n+1)!$ is tiny compared with $m$.
\end{solution}

\begin{solution}{ex:factorial}
Let $m=\lfloor n/2\rfloor\ge1$. By Bertrand there is a prime $p$ with $m<p\le2m\le n$; as $p\ge m+1>\frac n2$, we have $2p>n$. So $p\le n<2p$, and by Legendre $v_p(n!)=\left\lfloor\frac np\right\rfloor=1$, as $p^2\ge2p>n$. A $k$-th power has all valuations divisible by $k\ge2$.
\end{solution}

\begin{solution}{ex:erdosdiv}
(a) The number of $n\le N$ divisible by $p$ is $\lfloor N/p\rfloor\le N/p$; sum over $p_j$, $j>k$.

(b) Each such $n$ is $ab^2$ with $a$ a squarefree product of some of $p_1,\dots,p_k$ ($2^k$ choices) and $b\le\sqrt N$ ($\le\sqrt N$ choices).

(c) Every $n\le N$ falls into case (a) or (b), so $N<\frac N2+2^k\sqrt N$, that is $\sqrt N<2^{k+1}$. For $N=2^{2k+2}$ this reads $2^{k+1}<2^{k+1}$, a contradiction. So the tails of $\sum\frac1p$ never drop below $\frac12$, and the series diverges.
\end{solution}

\begin{solution}{ex:furstenberg}
(a) Unions of open sets are open by definition, and $\Z=0+1\Z$ is open (so is $\emptyset$). If $x\in U\cap V$ with $x+b\Z\subset U$ and $x+b'\Z\subset V$, then $x+bb'\Z\subset U\cap V$: finite intersections are open. A non-empty open set contains a progression, which is infinite.

(b) The complement of $a+b\Z$ is $\bigcup_{i=1}^{b-1}(a+i+b\Z)$, which is open.

(c) Every integer other than $\pm1$ has a prime factor, so $\Z\setminus\{-1,1\}=\bigcup_pp\Z$. With finitely many primes this is a finite union of closed sets, hence closed, so $\{-1,1\}$ is open. But it is non-empty and finite, contradicting (a).
\end{solution}
